\documentclass{article}
\usepackage{graphicx} % Required for inserting images
\usepackage{dsfont}
\usepackage{amsmath}
\usepackage{ dsfont }
\usepackage[all]{xy}
\usepackage{appendix}
\usepackage{stmaryrd} 
\DeclareMathOperator*{\argmin}{argmin}
\DeclareMathOperator*{\argmax}{argmax}
\usepackage{geometry}
\geometry{hmargin=2.5cm,vmargin=2.5cm}
\usepackage{amssymb}
\usepackage{caption}
\usepackage{subcaption}
\usepackage{amssymb}% http://ctan.org/pkg/amssymb
\usepackage{pifont}% http://ctan.org/pkg/pifont
\newcommand{\cmark}{\ding{51}}%
\newcommand{\xmark}{\ding{55}}%
\usepackage{xcolor}
\usepackage[french]{babel}
\usepackage{booktabs}
\usepackage{listings}
\usepackage{minipage-marginpar}

\newcommand*{\longDefiningEquals}{\stackrel{\text{def}}{=\joinrel=}}
\title{Rapport}
% \date{}

\usepackage[
   backend=biber,        % compilateur par défaut pour biblatex
   sorting=nyt,          % trier par nom, année, titre
   citestyle=numeric, % style de citation auteur-année
   bibstyle=numeric,  % style de bibliographie alphabétique
]{biblatex}
\addbibresource{ref.bib}

% \entrymodifiers={++[o][F-]}
\SelectTips{cm}{}
\begin{document}
\maketitle
\thispagestyle{empty} 
\newpage
\thispagestyle{empty} 
\tableofcontents
\newpage
\setcounter{page}{1}


\section{Définitions et hypothèses}\label{def}
On définit dans un premier temps l'ensemble des objets nécessaire à l'évaluation des contextes causaux contrefactuels. Une partie des définitions est directement extraite de l'article AMAAS \cite{sarmiento2025formalising}, d'autres ont été légérement modifiées et certaines ont été rajoutées spécifiquement pour l'étude.
\begin{itemize}
    \item $\mathds{F}$: ensemble de fluents qui décrivent l’environnement
    
    \item $\mathcal{F}$: ensemble des formules d'état, formules de logique classique construites sur $\mathds{F}$
    
    \item $\mathds{E}$: ensemble des variables qui décrivent les transitions, "ground events" qui modifient les fluents
    
    \item $\mathds{S} \subseteq 2^{\mathds{F}}$: ensemble de tous les états possibles du monde (exclusion des états incohérents)
    
    \item $\tau \subseteq \mathds{S} \times 2^{\mathds{E}} \times \mathds{S}$: ensemble des relations de transition. Pour chaque état indique les différents ensembles d'événements simultanés qui pourraient se produire ainsi que les états résultants
    
    \item $\langle\mathds{S}, \tau\rangle$: Labelled Transition System, graphe dirigé avec $\mathds{S}$ les nœuds du graphe et $\tau$ les arcs du graphe. nb: définition conventionnelle d'un graphe (nœuds + arcs), cette définition permet aussi de garder l'information d'états qui n'ont aucune transition entrante ou sortante. Soit $pre: \mathds{E}\to \mathcal{F}$, la fonction qui à chaque événement associe des préconditions on considère dans notre contexte que tous les LTS sont correct relativement à cette fonction.
    
    \item Politique $\pi$ : $\mathds{S} \to 2^{\mathds{E}}$ relativement à un LTS donné. A la différence de l'article AAMAS les événements ne dépendent que d'un état et non d'un couple (état,temps). On considère dans notre cadre que les états contiennent l'ensemble des informations nécessaires ce qui peut inclure des informations implicites sur le temps.

    \item $\Pi_\tau$ l'ensemble de toutes les politiques possibles relativement à un LTS $\langle\mathds{S},\tau\rangle$
    \item Politique correcte $\pi$ relative à un LTS: $\forall S \in \mathds{S} \quad \exists S'\in \mathds{S}, (S, \pi(S),S')\in \tau$

    


    \item Contexte causal: $\chi = (\pi, S_0,\tau)$ avec $\pi \in \Pi_\tau$ une politique, $S_0\in \mathds{S}$ un état de départ et $\tau$ un ensemble de relations de transitions décrivant un LTS. 
    
    \item Hypothèse cas déterministe: Soit un LTS défini par $\langle\mathds{S},\tau\rangle$ est dit déterministe si:
        \begin{equation*}
            \forall S,S'\in \mathds{S}, \forall E\in 2^{\mathds{E}}, \quad (S,E,S')\in \tau \implies \nexists S'' \in \mathds{S}, \quad (S,E,S'')\in \tau
        \end{equation*}
    \item $\mathds{A} \subseteq \mathds{S}$: ensemble des états terminaux d'un LTS $\langle\mathds{S},\tau\rangle$
        \begin{equation*}
            \mathds{A} = \{ S \mid S\in \mathds{S}, (S,\emptyset,S) \in \tau  \textit{ et }  \forall E\neq \emptyset,  \nexists S'\in \mathds{S}, (S,E,S')\in\tau\} \quad et \quad \forall S \notin \mathds{A}, \forall E \in 2^{\mathds{E}}, \quad (S,E,S)\notin \tau
        \end{equation*}
    \item $\mathds{T}$: ensemble des pas de temps
    
    \item $S^\chi: \mathds{T}\to \mathds{S}$ la fonction relative à un causal setting $\chi = (\pi,S_0,\tau)$ qui indique au pas de temps $t\in \mathds{T}$ l'état dans lequel on se trouve en suivant la politique $\pi$, dans $\tau$ depuis $S_0$.

    \item $E^\chi: \mathds{T}\to \mathds{S}$ la fonction relative à un causal setting $\chi = (\pi,S_0,\tau)$ qui indique au pas de temps $t\in \mathds{T}$ l'événement qui a lieu en suivant la politique $\pi$, dans $\tau$ depuis $S_0$
    
    \item $\sigma^\chi$: avec $\chi = (\pi,S_0,\tau)$, la trace complète (états, événements), obtenue dans le contexte causal $\chi$, de la forme: $[S_0,E_0,S_1,E_1,\dots,E_{n-1},S_n,E_n,S_n,E_n\dots] $. Une fois un état terminal atteint on observe une boucle qui amène à une trace infinie. On simplifie cette trace pour qu'elle soit de longueur finie. Soit $n$ tel que $S_n\in \mathds{A}$: la trace d'origine devient $\sigma^\chi=[S_0,E_0,S_1,E_1,\dots,E_{n-1},S_n]$. On notera $\sigma^{\chi}_{|state}$ la trace ne contenant que les états et, $\sigma^{\chi}_{|event}$ la trace ne contenant que les événements.
    

    \item $expert\_judge_\tau: \mathds{S} \times 2 ^{\mathds{E}}\to \lbrack 0, 1\rbrack$ fonction d'évaluation experte, définie partiellement sur l'ensemble des couples $(S,E)$ tel que $\exists S'\in\mathds{S},(S,E,S')\in \tau$. Cette fonction attribue à chaque paire une valeur comprise entre 0 et 1 relative à une norme (statistique ou prescriptive). Ces valeurs sont fixées de manière subjective.

\end{itemize}

\section{Exemple de transition d'un framework RAC vers LTS: PDDL}

Les LTS ne sont pas à proprement vu comme un langage pour le domaine du Raisonnement sur l'Action et le changement (RAC) mais, comme dit dans \cite{sarmiento_formalising_2025} leur utilisation permet d'unifier des raisonnement sur différents frameworks ainsi que de lever certaines ambiguités. Tous formalisme pouvant être représenté sous forme de LTS bénéficie ainsi des réflexions et observation de ce travail. Pour illustrer la possibilité de transition nous allons développer sur la traduction d'un problème PDDL vers LTS. La représentation LTS permet de travailler sur la comparaison de contexte causal contrefactuel mais on laisse le travail de recherche et de résolution de problème de plannification aux outils prééxistants. En effet la complexité exponentielle d'un LTS en fonction du nombre d'état et d'action n'est pas adapté à la recherche de politique.


\subsection{Rappel PDDL}

Un problème de planification est décrit en PDDL par deux types de fichiers. Le premier définit le domaine de planification il contient la déclaration des types, des actions (effets et préconditions) ainsi que les axiomes et les règles du domaine. Le second fichier définit le problème, il indique les objets, les constantes spécifiques ainsi que la situation de départ et l'objectif. Un exemple se trouve en annexe \ref{ex_pddl}


\subsection{Construction d'un LTS à partir d'un problème PDDL}
Pour construire un LTS à partir d'un couple de fichier domaine/problème on commence par instancier toutes les actions possibles à partir de leur définitions et des objets du problèmes, ce processus est explicité dans l'exemple ci-dessous.

\hfill
\begin{lstlisting}[language=Python, basicstyle=\small\ttfamily,
                   commentstyle=\color{gray}\itshape,
                   keywordstyle=\bfseries]
def instantiate(actions, objects):
    instanciated = []
    for action in actions:
        params = action.parameters
        for args in product(objects, repeat=len(params)):
            subst = dict(zip(params, args))
            name, precond,effect =  action.apply(subst)  
            instanciated.append((name, precond, effect))
    return instanciated
\end{lstlisting}
\hfill

Une fois l'ensemble de toutes les actions instanciées, avant de pouvoir cosntruire le LTS, on définit la fonction de satisfactions de préconditions. Pour un ensemble de préconditions et un état représenté par un ensemble de fluents, elle vérifie si les préconditions sont respectées dans l'état. Cette fonction gère les précondition négative, disjonctive et conjonctive.

\hfill
\begin{lstlisting}[language=Python, basicstyle=\small\ttfamily,
                   commentstyle=\color{gray}\itshape,
                   keywordstyle=\bfseries]
def statisfied(preconditions,etat):
    if preconditions conjonctives:
        return all(statisfied(p,etat) for p in preconditions)
    if preconditions negative:
        return not statisfied(preconditions,etat)
    if preconditions disjonctives:
        return any(statisfied(p,etat) for p in preconditions)
    return precondition in etat

\end{lstlisting}
\hfill

La seconde fonction qui est nécessaire avant celle de construction du graphe est la fonction d'application des effets positifs et négatifs d'une actions dans un etat.
\hfill

\begin{lstlisting}[language=Python, basicstyle=\small\ttfamily,
                   commentstyle=\color{gray}\itshape,
                   keywordstyle=\bfseries]
def apply_effects(effects,state):
    new_state = state
    for f in effects:
        if isinstance(f, Not):
            new_state.discard(eff) # effet negatif
        else:
            new_state.add(eff) #effet positif
    return new_state

\end{lstlisting}
\hfill

Avec ces trois fonctions on peut maintenant définir la fonction de construction du LTS associé au problème PDDL. Un état est défini par l'ensemble des prédicats qui sont vrais dans ce dernier. On utilise une pile pour progresser dans la construction du graphe, à chaque nouvel état on observe parmis toutes les actions possibles lesquelles sont applicables en fonction de leur préconditions. Si une action est applicable on ajoute une nouvelle arête qui a pour état d'arrivée soit un état préexistant soit un nouvel état que l'on ajoute à la pile. Les états finaux sont les états dans lesquelles les conditions de l'objectif sont respectés. Pour les distinguer des autres états une arête réflexive est ajoutée. 

\hfill
\begin{lstlisting}[language=Python, basicstyle=\small\ttfamily,
                   commentstyle=\color{gray}\itshape,
                   keywordstyle=\bfseries]
def build_graph(actions, init_state, goals):
    state = init_state
    G = graph
    G.add_node(state)
    queue = [state]
    visited = {}
    while queue:
        current_node = queue.pop(0)
        if current_node in visited:
            continue
        visited.add(current_node)

        for preconditions, effects in actions:
            if satisfied(preconditions, current_node):
                new_node = apply_effect(effects, current_node)

            exist = new_node in G.nodes
            if not exist:
                G.add_node(new_node)
            else:
                continue
            
            G.add_edge(current_node, new_node)

            if new_node not in visited:
                queue.append(visited)
    
    for node in G.nodes:
        if satisfied(goals, node):
            G.add_edge(node,node, action = '$\\varnothing')

\end{lstlisting}
\hfill

On obtient pour l'exemple \ref{ex_pddl} le graphe suivant:

\begin{figure}[h!]
    \[
        \xymatrix @R=0.1em @C=5em {
            *++[o][F-]{S0} \ar@/_/[r]_{unstack\ a\ b} & 
            *++[o][F-]{S1} \ar@/_/[l]_{stack\ a\ b} \ar@/_/[r]_{put\ down\ a} &
            *++[o][F-]{S2} \ar@/_/[l]_{pick\ up\ a} \ar@/_/[r]_{pick\ up\ b} \ar@(r,u)[]_{\varnothing} &
            *++[o][F-]{S3} \ar@/_/[l]_{put\ down\ b} \ar@/_/[r]_{stack\ b\ a} &
            *++[o][F-]{S4} \ar@/_/[l]_{unstack\ b} \ar@(r,u)[]_{\varnothing}
        }
    \]
    \centering
    \captionsetup{width=0.8\textwidth}
    \caption{LTS obtenu pour le problème simple de block world. Avec S0 =\{'clear a', 'on a b', 'on table b', 'arm-empty'\}, S1 =\{'on table b', 'clear b', 'holding a'\}, S2 =\{'on table b','arm empty'', 'clear b','clear a', 'on table a'\}, S3 =\{'holding b','on table a', 'clear a'\}, S4 =\{'on b a', 'on table a', 'arm empty', 'clear b'\}}
    \label{fig:exemple graphe BW}
\end{figure}

\subsection{Remarques et limites}

Les LTS obtenus ne représentent pas l'ensemble des états possibles du monde, seul les états accessibles depuis l'état initial, respectant les fonctions de précondition y figurent. Pour enrichir ces graphes il serait potentiellement intéressant de construire pour un même domaine et objectif,  différents graphes selon différents états de départ et les fusionner afin de pouvoir travailler notamment sur des contextes contrefactuels avec des états de départ qui diffèrent. Cependant les LTS correspondent à une sémantique et non pas à de vrais outils de travail, leur utilité pour des applications concrètes reste limitée et ils seront surtout utilisés de manière formelle pour raisonner. Pour des applications à des exemples simples ils permettent une meilleur visualisation, on considère donc suffisant un parser simple comme présenté ci-dessus. L'idée de cette section étant principalement d'expliciter la faisabilité d'une transition d'un framework RAC vers des LTS, de possibles améliorations sont laissées en suspends car ne représentant pas le coeur du travail. On peut par exemple noter que ce parser ne gère que des problème PDDL contenant uniquement le requirement STRIPS. 

Il est possible de continuer de tirer parti de PDDL et de ses solveurs pour extraire de manière automatique des politiques qui nous permettent par la suite de travailler sur des contextes causaux. Pour construire la politique on avance dans le graphe depuis l'état de départ en attribuant à chaque état dans lequel on se trouve l'action associée dans le plan. Ces politiques partielles peuvent ensuite être complétées aléatoirement ou selon une heuristique.



\section{Évaluation d'un contexte causal contrefactuel}


Pour évaluer un contexte causal contrefactuel relativement à un contexte causal observé on procède en deux temps. Dans un premier temps on attribue un coût au contexte causal, ce dernier permet d'évaluer la proximité à l'observation, ainsi que la plausibilité du contexte. Le coût permet aussi d'agir comme un filtre à priori sur les contrefactuels non valides relativement au LTS. Une fois le coût calculé on applique un critère de validité (au sens XAI), qui permet de vérifier que les sorties obtenues, en appliquant respectivement les deux politiques, à partir de l'état de départ dans le LTS, sont différentes. Ce critère agit comme un filtre à posteriori sur les contextes causaux contrefactuels.



\subsection{Fonction de coût}

La fonction de coût doit permettre un compromis entre l'évaluation:
\begin{itemize}
    \item  de la distance du contexte causal contrefactuel par rapport au causal setting  observé. L'intuition derrière ce terme corresponds à l'idée qu'un bon contrefactuel doit être relativement proche du contexte causal observé pour être pertinent.
    \item de la plausibilité du contexte causal. La plausibilité ou encore le réalisme corresponds ici à un respect relatif d'une norme (statistique ou prescriptive). 

\end{itemize}

La fonction de coût intègre deux termes supplémentaires qui permettent d'exclure les contrefactuels non valides relativement au LTS (politiques non correctes), ainsi que les politiques qui ne divergent de la politique observée que sur des états non présents dans la trace d'états observée.

\begin{equation}\label{cout}
    \textit{Coût}_\tau(\chi,\chi')= \lambda_1  d\_obs(\sigma^\chi_{|state},\sigma^{\chi'}_{|state}) + \lambda_2 v(\chi', expert\_{judge}_\tau) + divergence(\chi,\chi') + correct(\chi') 
\end{equation}

Avec:
\begin{itemize}
    \item $\chi= (\pi,S_0,\tau)$ le contexte causal observé et $\chi' = (\pi',S_0',\tau)$ le contexte causal contrefactuel que l'on évalue. 
    \item $d\_obs$ attribue un coût relatif à la distance entre le contexte causal contrefactuel par rapport à celui observé. On utilise la distance de Levenshtein, avec variation des poids en fonction de l'intervention effectuée (suppression, insertion, substitution).
    \item $v$ la fonction de vraisemblance du contexte causal contrefactuel par rapport à un jugement expert.
    \item La fonction $divergence$ permet d'exclure les contextes causaux contrefactuels si leur politique associée ne diverge de celle observée que dans des états non présents dans la trace d'états du contexte causal observé. 
    \begin{equation}
        divergence(\chi,\chi') = \begin{cases}
             & 0 \quad \text{si} \quad \exists t<len(\sigma^{\chi}_{|state}) \quad \pi(S^\chi(t))\neq \pi'(S^{\chi}(t))
             \\
             &+\infty \quad \text{sinon}
        \end{cases}
    \end{equation}
    Voir exemple \ref{ex_div}

    \item La fonction $correct$ permet d'exclure les contextes causaux contrefactuels dont la politique n'est pas correcte:
    \begin{equation}
        correct(\chi') = \begin{cases}
            & 0 \quad \textit{si $\pi'$ correcte relativement au LTS associé à $\tau$} \\
            & +\infty \quad sinon
        \end{cases}
    \end{equation}
    Voir exemple \ref{ex_corr}.

\end{itemize}
 % \subsection{Distance à l'observation}
 \subsection{Vraisemblance du contrefactuel}
 Pour évaluer la vraisemblance (fonction $v$) du contexte causal contrefactuel on se base sur l'évaluation extérieure d'un juge expert, comme définie dans la partie \ref{def}. Cette évaluation permet d'attribuer un score compris entre 0 et 1 à chaque couple état/événement d'un LTS. Un score de 1 corresponds à un respect total de la norme, un score de 0 à un non respect total de la norme. 
 À partir des scores on évalue la vraisemblance du contrefactuel soit en l'exploitant directement, soit en construisant à partir de ces derniers une politique optimale.
 
 \subsubsection{Juge expert}
Différentes méthodes sont applicables pour attribuer les scores. On peut contraindre, par exemple, pour un état donné, à avoir une distribution de probabilité ou de possibilité sur l'ensemble des événements possibles depuis cet état.
 \begin{enumerate}
    \item Distribution de probabilité: L'expert donne une distribution de probabilité sur les arcs sortant d'un état, deux possibilités de construction:
        \begin{itemize}
            \item À partir d'un classement avec softmax sur l'inverse des positions: simple mais on perds en expressivité, les écarts entre les événements sont fixes relativement au nombre d'arcs sortants du sommet qui représente un état. De plus cela demande d'avoir un ordre sur les événements ce qui n'est pas toujours souhaitable si on a des événements normativement équivalents.
            \item L'expert donne directement une distribution de probabilités, cette méthode a l'avantage de limiter la perte d'expressivité et d'autoriser des scores égaux pour deux événements sortant d'un même état.
            \item Une distribution de probabilité peut fortement pénaliser des événements issus d'un état avec de nombreux arcs sortant comme on l'observe sur le LTS ci-dessous:
        \end{itemize}

        \begin{figure}[h!]\label{graph_prob}
            \centering
            \[
                \xymatrix @R=0.1em@C=5em {
                    & & *++[o][F-]{D} \ar[r]^{\dots} &
                \\ & & *++[o][F-]{E} \ar[r]^{\dots}&
                \\ & *++[o][F-]{B} \ar[ruu]^{0.25}\ar[ru]^{0.25}\ar[r]^{0.25}\ar[rd]^{0.25}& *++[o][F-]{B}&
                \\ & & *++[o][F-]{G} \ar[r]^{\dots}&
                \\  *++[o][F-]{A} \ar[ruu]^{0.5} \ar[rd]^{0.5}& & &
                \\ & *++[o][F-]{C}\ar[r]^{0.5} \ar[rd]^{0.5}&*++[o][F-]{H} \ar[r]^{\dots} &
                \\ & & *++[o][F-]{I} \ar[r]^{\dots}& 
                }            
            \]
            \captionsetup{width = 0.6\textwidth}
            \caption{Exemple de LTS contraint à des distributions de probabilités. Pour les états B et C chaque événement sortant est normativement équivalent aux autres. Les arcs sortant de l'état B sont pénalisés par leur nombre.}
            \label{fig:placeholder}
        \end{figure}

    \item Distribution de possibilité: l'expert attribue un distribution de possibilité sur l'ensemble des arcs sortant d'un état donné. On note qu'une distribution de possibilité impose que le score maximal parmi tous les scores des événements sortant d'un état soit de 1.
    \item Valeur libres: l'expert juge indépendamment chaque événement possible en attribuant une valeur entre 0 et 1. Avantage: pas de contrainte sur la somme des scores ou encore sur leur valeur maximale. Désavantage: interprétation et choix des valeurs potentiellement plus complexe.  
        
\end{enumerate}
 \subsubsection{Construction d'une politique optimale à partir de l'expert}
 Pour évaluer la vraisemblance d'un contexte causal contrefactuel à partir d'une évaluation experte, on peut premièrement, construire une politique optimale vis à vis de l'évaluation puis dans un second temps comparer nos deux politiques (optimale et contrefactuelle) . La construction de politiques optimales peut être contrainte selon l'application, en fonction d'une heuristique, en prenant en compte le coût de calcul de la construction, etc...Les politiques optimales dépendent uniquement d'un LTS et de l'évaluation d'un juge expert, elles sont indépendantes des contextes causaux contrefactuels qui sont évalués.
 \begin{enumerate}
     \item Politique optimale sans planification: pour chaque état cette politique indique l'événement qui maximise l'évaluation de l'expert à un horizon temporelle d'un seul pas de temps.  La fonction $\pi_\tau^{opt}\in \Pi_\tau$ est définie par:
     \begin{equation}\label{pi_instan}
            \forall S\in \mathds{S} \quad \pi^{opt}_\tau(S) = \underset{E}{\argmax} (expert\_{judge}_\tau(S,E))
     \end{equation}
     Cette politique a pour avantage sa simplicité de construction, d'interprétation ainsi que son coût de calcul. Dans certains cas elle possède tout de même certaines limitations comme on l'observe dans l'exemple ci-dessous:

     \begin{figure}[h!]
            \[
                \xymatrix @R=0.1em{
                    & &  *++[o][F-]{S_3}\ar@(r,u)[]_{\emptyset} &
                \\  & *++[o][F-]{S_1} \ar@[red][ru]_{a: 0.2} \ar[rd]_{b: 0.2}&     
                \\  & & *++[o][F-]{S_4}\ar@(r,u)[]_{\emptyset} &
                \\  *++[o][F-]{S_0} \ar@[red][ruu]_{a: 0.7} \ar@[green][rdd]_{b: 0.4}& &   
                \\  & & *++[o][F-]{S_5} \ar@(r,u)[]_{\emptyset} &   
                \\  & *++[o][F-]{S_2} \ar[ru]_{a: 0.8} \ar@[green][rd]_{b: 0.9}&     
                \\  & & *++[o][F-]{S_6}\ar@(r,u)[]_{\emptyset} &
                }
            \]
            \captionsetup{width=0.8\textwidth}
         \centering
            \caption{Exemple de LTS évalué par un expert qui représente un cas limite pour une politique sans planification. Étiquetage des arcs du graphe de la forme "événement:score"}
         \label{fig:instant_cas_lim}
     \end{figure}
     La politique optimale sans planification associée au LTS de la Figure \ref{fig:instant_cas_lim} corresponds à:
     \begin{equation}
         \pi_\tau^{opt1}(S_0)=a,\quad \pi_\tau^{opt1}(S_1)=a,\quad \pi_\tau^{opt1}(S_2)=b,\quad \forall i\in[3,4,5,6]\quad  \pi^{opt1}_\tau(S_i)=\emptyset
     \end{equation}
    Si on fixe l'état de départ à $S_0$ on obtient, en suivant la politique optimale, la trace: $[S_0, a, S_1,a,S_3]$ qui corresponds au chemin rouge. La moyenne arithmétique des scores des événements de cette trace vaut 0.45 (le  choix de la moyenne arithmétique est discutable). Cependant on peut considérer dans ce cas que la bonne décision à prendre, à partir de l'état $S_0$, en connaissant le LTS et les scores qui y sont associés est l'événement $b$. En effet on obtient la trace $[S_0,b,S_2,b,S_6]$, correspondant au chemin vert et qui a une moyenne arithmétique sur les scores des événements de la trace de 0.65. 
    
    De cet exemple on tire deux observations. La politique optimale sans planification ne tire pas partie de la connaissance totale du LTS et des scores associés aux événements. Deuxièmement on note que la politique optimale n'est pas unique. En effet, un choix arbitraire a été effectué en $S_1$ pour le choix de l'événement $a$. Dans les cas où plusieurs événements maximisent le score on peut appliquer une seconde heuristique ou choisir aléatoirement.
     
     \item Politique optimale avec planification: on peut généraliser la définition \ref{pi_instan} à un horizon temporel $n$. Pour cela, dans chaque état on applique un algorithme du plus court chemin dans un graphe pondéré (par exemple A*) et on sélectionne le premier événement de ce chemin. Voir \ref{astar} pour plus de détail sur la construction de cette politique. Ce type de politique optimale permet de prendre en compte la connaissance totale du LTS et en fonction de l'algorithme de plus court chemin on peut limiter l'augmentation du coût de construction de la politique. Une telle politique permet de mieux gérer les cas du même type que la Figure \ref{fig:instant_cas_lim}. On obtient, par exemple pour le LTS de la figure \ref{fig:instant_cas_lim} une politique de la forme:
     \begin{equation}
         \pi_\tau^{opt2}(S_0)=b,\quad \pi_\tau^{opt2}(S_1)=a,\quad \pi_\tau^{opt2}(S_2)=b,\quad \forall i\in[3,4,5,6]\quad  \pi^{opt2}_\tau(S_i)=\emptyset
     \end{equation}
     Comme pour les politiques optimales sans planification, il n'y a pas unicité du résultat. Un traitement spécifique est nécessaire pour les états qui possèdent plusieurs événements qui maximise le score.
 \end{enumerate}

\subsubsection{Évaluation de la vraisemblance à partir de la politique optimale}

À partir d'une politique optimale donnée on peut définir un coût de vraisemblance entre cette dernière et la politique associé à un contexte causal contrefactuel. On définit une fonction de distance entre deux politiques sur une trace d'état: 
\begin{equation}
dist\_politiques(\pi,\pi',\sigma_{|state}) =\frac{ \#\{\pi(S) \neq \pi'(S)| S\in \sigma_{|state}\}}{\#\{S\in \sigma_{|state}\}}    
\end{equation}

On peut utiliser cette distance pour évaluer la vraisemblance entre une politique associée à un contexte causal contrefactuel $\chi' = (\pi',S_0,\tau)$ et une politique optimale $\pi^{opt}_\tau$ construite à partir d'une évaluation experte $expert\_judge_\tau$:
\begin{equation}
    v(\chi',expert\_judge_\tau)= dist\_politiques(\pi',\pi^{opt}_\tau, \sigma^{\chi'}_{|state})
\end{equation}
On compte le nombre de fois où la politique contrefactuelle n'a pas effectué un choix optimal sur le chemin (la trace) obtenue dans son contexte causale.
\newline

Exemple: Soit le LTS de la Figure \ref{fig:instant_cas_lim} $\langle\mathds{S},\tau\rangle$ on définit un contexte causal $\chi'=(\pi',S_0,\tau)$ avec:
$$\pi'(S_0)= b,\quad \pi'(S_1) = b,\quad \pi'(S_2)=a,\quad \forall i\in[3,4,5,6]\quad  \pi'(S_i)=\emptyset$$
On obtient la trace $\sigma^{\chi'}=[S_0,b,S_2,a,S_5]$ et donc $\sigma_{|state}^{\chi'}=[S_0,S_2,S_5]$
\begin{itemize}
    \item Coût de vraisemblance avec $\pi^{opt1}_\tau$:
    \begin{equation*}
        dist\_politiques(\pi',\pi^{opt1}_\tau, \sigma^{\chi'}_{|state}) = \frac{\#{\{S_0,S_2\}}}{\#\{S_0,S_2,S_5\}} = \frac{2}{3}
    \end{equation*}
    
    \item Coût de vraisemblance avec $\pi^{opt2}_\tau$:
    \begin{equation*}
        dist\_politiques(\pi',\pi^{opt2}_\tau, \sigma^{\chi'}_{|state}) = \frac{\#{\{S_2\}}}{\#\{S_0,S_2,S_5\}} = \frac{1}{3}
    \end{equation*}
\end{itemize}
On observe que le résultat du coût de vraisemblance dépends de la politique optimale même si l'évaluation experte sur lesquelles $\pi^{opt1}_\tau$ et $\pi^{opt2}_\tau$ sont construites est la même.

\subsubsection{Exploitation directe de l'expert}
Une autre approche possible pour évaluer la vraisemblance d'un contexte causal contrefactuel est de directement exploiter l'évaluation experte, en agrégeant les scores des événements associés à la trace du contexte causal contrefactuel. Pour cela différents opérateurs d'agrégation peuvent être utilisés en fonction de l'application.



Soit $Agg$ un opérateur d'agrégation, un LTS $\langle \mathds{S},\tau \rangle$, un contexte causal contrefactuel $\chi'= (\pi',S_0,\tau)$ avec $\sigma^{\chi'}= [S^{\chi'}(0),E^{\chi'}(0),\dots,E^{\chi'}(n-1),S^{\chi'}(n)]$ et $expert\_judge_\tau$ une évaluation experte sur le LTS, la fonction $v$ de l'équation \ref{cout} est de la forme:

\begin{equation}
    v(\chi',expert\_judge_\tau)= 1 - Agg(s_0, \dots, s_{n-1})) \quad \forall i < len(\sigma^{\chi'}_{|state}) \quad s_i= expert\_judge(S^{\chi'}(i),E^{\chi'}(i))
\end{equation}\label{exploit_direct}



Différentes propriétés peuvent être désirables pour les opérateurs d'agrégation: une invariance à la longueur de la liste des scores (implicitement du chemin), une pénalisation en fonction du temps, des possibilités de compensation, etc...


 \paragraph{Opérateur moyenne géométrique} 
 Soit $\epsilon>0$ qui permet de ne pas avoir à restreindre l'évaluation à des valeurs strictement supérieur à 0.
    \begin{equation}
            agg_{geo}(s_0,s_1,\dots,s_{n-1}) =\sqrt[n]{\prod_{i=0}^{n-1} \max (\epsilon,s_i)} 
            % agg_{geo}(s_0,s_1,\dots,s_{n-1}) =\exp\left( \frac{1}{N} \sum_{i<n} \log(\max(s_i,\epsilon ))\right) 
    \end{equation}
    Cet opérateur permet une plus grande expressivité que la moyenne arithmétique. Il est plus sensible à une variance importante. Même si il reste un opérateur de compensations, pour "rattraper" un score faible, la compensation doit être plus forte qu'avec une moyenne arithmétique. Par exemple:
    \begin{itemize}
        \item $1 - agg_{geo}([0.5,0.5,0.5]) = 0.5$
        \item$ 1 - agg_{geo}([0.5,0.3,0.7]) = 0.53$ la moyenne arithmétique est bien de 0.5 mais, l'opérateur pénalise le faible score de 0.3
        \item $1 - agg_{geo}([0.5,0.3,0.8]) = 0.5$ la moyenne arithmétique est de 0.53
        
    \end{itemize}

\begin{figure}[h!]
    \centering
    \begin{subfigure}[b]{0.45\textwidth}
        \centering
        \[
        \xymatrix @R=0.1em@C=5em {
         & *++[o][F-]{B} \ar[r]^{b:0.4} & *++[o][F-]{C} \\
          *++[o][F-]{A} \ar[ru]^{a:0.4} \ar[rd]_{c:0.1} \\
        & *++[o][F-]{D} \ar[r]_{d:0.9} & *++[o][F-]{E}
        }
        \]
        \caption{LTS}
        \label{fig:lts_geo}
    \end{subfigure}
    \hfill
    \begin{subfigure}[b]{0.45\textwidth}
        \centering
        \begin{tabular}{ccc}
            chemin & 1 - moyenne & 1 - $agg_{geo}$ \\
            \hline
            ABC & 0.6 & \textbf{0.6} \\
            ADE & \textbf{0.5 }& 0.7 \\
        \end{tabular}
        \caption{Coûts}
        \label{fig:Coûts}
    \end{subfigure}
    \caption{Exemple opérateur géométrique}
    \label{fig:geo}
\end{figure}
On observe sur la figure \ref{fig:geo} que le coût associée au chemin [ABC] est plus faible que celui du chemin [ADE] même si ce dernier possède une moyenne arithmétique plus élevée. En effet, la variance plus importante de [ADE] est pénalisée en utilisant l'opérateur géométrique. On peut imaginer des situations simple où l'on préfère faire une action neutre comme mettre du sirop dans un verre d'eau plutôt que de mettre du poison puis un antidote avant de le servir (en utilisant une norme relative au bien et au mal).
    
\paragraph{Opérateur gamma (temporel)}
Soit $\gamma\in ]0,1]$, $\Gamma = \frac{1-\gamma^{N}}{1-\gamma}$ ($\Gamma =N $ si $\gamma=1$)
    \begin{equation}
        agg_{\gamma}(s_0,s_1,\dots,s_{n-1}) = \frac{ \sum_{i=0}^{N-1}\gamma^is_i}{\Gamma}
    \end{equation}
    Cet opérateur permet d'accorder une grande importance aux événements au début d'une trace puis de manière exponentiellement décroissante à ceux qui arrivent à un temps avancé. On peut facilement définir de manière symétrique l'agrégateur $agg_{\gamma \text{ }late}$:
    \begin{equation}
        agg_{\gamma \text{ }late}(s_0,s_1,\dots,s_{n-1}) = \frac{ \sum_{i=0}^{N-1}\gamma^{N-1-i}s_i}{\Gamma}        
    \end{equation}
    
    Ce dernier accorde plus d'importance aux scores en fin de traces qu'à ceux du début (voir Figure \ref{fig:gamma}). La valeur du $\gamma$ possède un rôle essentiel pour cet opérateur et doit être fixée relativement à la longueur de la liste plus la liste des scores est longue plus $\gamma$ doit être proche de 1 et inversement pour des listes courtes, il est préférable de fixer $\gamma$ à des valeurs proches de 0.
    
    
    % \begin{figure}[h!]
    %     \centering
    %     \includegraphics[width=0.7\linewidth]{gammas_opp.png}
    %     \caption{Évolution des scores des agrégateurs gamma et gamma late en fonction de la position d'un événement de faible score dans la liste des scores. La liste des scores possède 100 éléments qui ont tous une valeur de 1. L'élément de faible valeur que l'on fait progressivement avancer dans la liste a une valeur de 0.1, $\gamma=0.9$}
    %     \label{fig:placeholder}
    % \end{figure}
    On peut utiliser l'opérateur gamma si l'on souhaite pénaliser des événements avec un faible score en début de trace. Par exemple, si on utilise une norme statistique, on souhaite exclure les chemins qui dès les premiers temps possèdent un événement statistiquement peu probable même si après tous les événements qui en découlent ont une forte possibilité de se produire. 


    \begin{figure}[h!]
        \centering
        \begin{subfigure}[b]{0.45\textwidth}
            \centering
            \[
            \xymatrix @R=0.1em@C=5em {
             & *++[o][F-]{B} \ar[r]^{b:0.9} & *++[o][F-]{C} \ar[r]^{c:0.9} & *++[o][F-]{D}  \\
              *++[o][F-]{A} \ar[ru]^{a:0.1} \ar[rd]_{d:0.9} \\
            & *++[o][F-]{E} \ar[r]_{e:0.9} & *++[o][F-]{F} \ar[r]_{f:0.1}  & *++[o][F-]{G}
            }
            \]
            \caption{LTS}
            \label{fig:lts_gamma}
        \end{subfigure}
        \hfill
        \begin{subfigure}[b]{0.45\textwidth}
            \centering
            \begin{tabular}{cccc}
                chemin & 1 - moyenne & 1 - $agg_{\gamma}$ & 1- $agg_{\gamma late}$  \\
                \hline
                ABCD & \textbf{ 0.3} & 0.42 &\textbf{ 0.21} \\
                AEFG & \textbf{0.3} & \textbf{0.21} & 0.42\\
            \end{tabular}
            \caption{Coûts ($\gamma =0.7$)}
            \label{fig:Coûts}
        \end{subfigure}
        \caption{Exemple opérateurs gamma}
        \label{fig:gamma}
    \end{figure}


    
    Les opérateurs gamma permettent de prendre en compte le temps dans l'agrégation des scores, ils ne sont pas invariants à la longueur de la trace et ne permettent pas une réelle compensation entre des scores si ils sont éloignés temporellement.

\paragraph{Opérateur seuil} soit $\lambda\in[0,1]$
\begin{equation}
    agg_{seuil}(s_0,\dots,s_{n-1}) = \begin{cases}
        0 & si \quad \exists i<n \textit{ tel que } s_i < \lambda
        \\ 1 & sinon
    \end{cases}
\end{equation}



L'opérateur de seuil permet soit de pénaliser complètement un contrefactuel si un événement possède un score inférieur au seuil fixé soit de ne pas pénaliser si un tel événement n'apparaît pas dans la trace voir Figure \ref{fig:seuil}. 

    \begin{figure}[h!]
        \centering
        \begin{subfigure}[b]{0.45\textwidth}
            \centering
            \[
            \xymatrix @R=0.1em@C=5em {
             & *++[o][F-]{B} \ar[r]^{b:0.3} & *++[o][F-]{C} \ar[r]^{c:0.3} & *++[o][F-]{D}  \\
              *++[o][F-]{A} \ar[ru]^{a:0.3} \ar[rd]_{d:0.9} \\
            & *++[o][F-]{E} \ar[r]_{e:0.2} & *++[o][F-]{F} \ar[r]_{f:0.9}  & *++[o][F-]{G}
            }
            \]
            \caption{LTS}
            \label{fig:lts_gamma}
        \end{subfigure}
        \hfill
        \begin{subfigure}[b]{0.45\textwidth}
            \centering
            \begin{tabular}{ccc}
                chemin & 1 - moyenne & 1 - $agg_{seuil}$   \\
                \hline
                ABCD & 0.6 & \textbf{0}  \\
                AEFG & \textbf{0.33} & 1 \\
            \end{tabular}
            \caption{Coûts ($\lambda =0.25$)}
            \label{fig:Coûts}
        \end{subfigure}
        \caption{Exemple opérateurs de seuil}
        \label{fig:seuil}
    \end{figure}

Cet opérateur peut présenter un intérêt dans des cas critiques où l'on se refuse de passer sous un seuil de score de normativité quel que soit le reste des événements. 

\subsubsection{Sélection des événements à évaluer}

On peut vouloir agréger l'ensemble des scores d'un chemin causal (comme fait jusqu'alors dans les exemples). Mais comme le travail porte in fine sur une évaluation de contrefactuel on peut vouloir se concentrer uniquement sur les transitions où une alternative existe, c'est à dire les états avec strictement plus d'une transition sortante. 



\begin{figure}[h!]
    \centering
    \begin{subfigure}[b]{0.45\textwidth}
        \centering
        \[
            \xymatrix @R=0.1em{
                &  *++[o][F-]{B} \ar[r]^{d:1} \ar@[blue]@{-}@<2pt> [r]& *++[o][F-]{D} \ar[r]^{g:0.9} \ar@[blue]@{-}@<2pt>[r] &  *++[o][F-]{G} \ar@[blue]@{-}@<2pt> [r] \ar[r]^{i:1} & *++[o][F-]{I}
                \\ *++[o][F-]{A} \ar[ru]^{b:0.2} \ar[rdd]^{c:0.8}  \ar@[red]@{-}@< 2pt> [rdd] \ar@[green]@{-}@< -2pt> [rdd] \ar@[blue]@{-}@<2pt> [ru] & & & 
                \\ & & *++[o][F-]{E} \ar[r]^{h:0.3} \ar@[red]@{-}@<2pt> [r] & *++[o][F-]{H} &
                \\ & *++[o][F-]{C} \ar[ru]^{e:0.8} \ar@[red]@{-}@<2pt> [ru] \ar[rd]^{f:0.3} \ar@[green]@{-}@< -2pt> [rd]  & & &
                \\ & & *++[o][F-]{F}
            }
        \]
        \caption{LTS}
        \label{fig:lts_gamma}
    \end{subfigure}
    \hfill
    \begin{subfigure}[b]{0.45\textwidth}
        \centering
        \begin{tabular}{ccc}
            chemin & 1 - $agg_{geo}$ & événements comptés \\
            \hline
            bleu  & 0.\textbf{35} & $b,d,g,i$ \\
            rouge & 0.42 & $c,e,h$ \\
            vert  & 0.51 & $c,f$\\
        \end{tabular}
        \caption{Coûts globaux ($\gamma = 0.7$)}
        \label{tab:couts_global}

        \vspace{1em}

        \begin{tabular}{ccc}
            chemin & 1 - $agg_{geo}$  & événements comptés \\
            \hline
            bleu  & 0.8 & $b$ \\
            rouge & \textbf{0.2} & $c,e$ \\
            vert  & 0.51 & $c,f$ \\
        \end{tabular}
        \caption{Coûts sur les choix ($\gamma = 0.7$)}
        \label{tab:couts_choix}
    \end{subfigure}
    \caption{Évaluation sur les choix}
    \label{fig:choix}
\end{figure}

 Comme on l'observe sur la Figure \ref{fig:choix}, considérer seulement les événements avec une alternative modifie complètement les résultats, on évalue la qualité des choix pris durant un scénario et non plus le scénario dans sa globalité. Lorsque l'on évalue la normativité totale c'est le chemin bleu qui minimise le coût, hors ce scénario pour l'unique choix possible qu'il a, a effectué le pire choix. Contrairement à ce dernier, le chemin rouge sur les deux événements qui présentent une alternative prends les événements qui maximise le score. En réduisant les événements à considérer c'est le chemin rouge qui minimise le coût.

\subsection{Extension de la fonction de coût: distance entre événements}

La distance à l'observation est essentielle pour évaluer les contextes causaux contrefactuels, actuellement le terme $d\_obs$ est l'unique terme qui nous permet d'évaluer cette distance. Ce terme est calculé comme la distance de Levenshtein sur la trace d'état d'un contexte causal, hors dans des multigraphes cette distance peut être insuffisante. 

\subsection{Critère de validité}
Le critère permet de filtrer les contrefactuels non valides. Un contexte causal contrefactuel et considéré comme valide par rapport au contexte causal observé si l'état terminal de la trace du contrefactuel est différent de celui de l'observation et que son coût est inférieur à un seuil fixé.

Soit $\chi= (\pi,S_0,\tau)$ le contexte causal observé avec $\sigma_{|state}^\chi = [S^\chi(0),\dots,S^\chi(n)]$ et $\chi'= (\pi',S_0',\tau)$ le contexte causal contrefactuel que l'on évalue, avec $\sigma_{|state}^{\chi'} = [S^{\chi'}(0),\dots,S^{\chi'}(m)]$ et $\delta \in \mathds{R}^+$, $\chi'$ est valide si:
\begin{equation}
    S^\chi(n) \neq S^{\chi'}(m) \quad \text{et} \quad \textit{Coût}_\tau(\chi,\chi')<\delta
\end{equation}

On peut aussi envisage d'affiner le critère de validité en rajoutant des seuils propres à $v$ et $d\_obs$.

\section{Exemple Poison}

\subsection{Définitions du LTS}


$\mathds{F} = \{VerreOK, \textit{VerreEmpoisonné},VerreVide,VerreAntidote,Mort,Vivant\}\\ \quad \mathds{E} = \{empoisonner,antidote,boire,vider\}\\
\mathds{T}= \mathds{N} \\
pre(boire) = VO \vee VA \vee VE \vee VV \quad pre(antidote) = VE \quad pre(empoisonner)= VO \quad pre(vider) = VO\vee VE$\\
\[
\mathbb{S} = \left\{
\begin{aligned}
VO     &= VO     \wedge \neg VE \wedge \neg VV \wedge \neg VA \wedge \neg \text{Mort} \wedge \neg \text{Vivant}, \\
VE     &= \neg VO \wedge VE     \wedge \neg VV \wedge \neg VA \wedge \neg \text{Mort} \wedge \neg \text{Vivant}, \\
VV     &= \neg VO \wedge \neg VE \wedge VV     \wedge \neg VA \wedge \neg \text{Mort} \wedge \neg \text{Vivant}, \\
VA     &= \neg VO \wedge \neg VE \wedge \neg VV \wedge VA     \wedge \neg \text{Mort} \wedge \neg \text{Vivant}, \\
\text{Mort}   &= \neg VO \wedge \neg VE \wedge \neg VV \wedge \neg VA \wedge \text{Mort}   \wedge \neg \text{Vivant}, \\
\text{Vivant} &= \neg VO \wedge \neg VE \wedge \neg VV \wedge \neg VA \wedge \neg \text{Mort} \wedge \text{Vivant}
\end{aligned}
\right\}
\]

$\mathds{A} = \{ Vivant,Mort\}$
\[
\tau = \left\{
\begin{aligned}
(VE,vider,VV),(VE,boire,Mort),(VE,antidote,VA),(Mort,\emptyset,Mort)\\
(VV,\emptyset,Mort),(VA,boire,Vivant),(Vivant,\emptyset,Vivant),(VO, boire,vivant)\\
(VO, boire,vivant),(VO,empoisonner,VE),(VO,vider,VV),\\
\end{aligned}
\right\}
\]

\begin{center}
\[
expert\_judge_\tau= \left\{
\begin{aligned}
VO &: \{boire: 1,\ vider: 0.5,\ emp: 0.1\}, \\
VE &: \{antidote: 1,\ boire: 0.1,\ vider: 0.5\}, \\
VV &: \{\emptyset: 0.5\}, \\
VA &: \{boire: 1\}, \\
\text{Mort} &: \{\emptyset: 1\}, \\
\text{Vivant} &: \{\emptyset: 1\}
\end{aligned}
\right\}
\]
\end{center}

\begin{align}
    \scalebox{1.2}{
        \xymatrix@=2em@C=5em{
            &    *++[o][F-]{Vivant} \ar @(r,u) []_{\emptyset}^1& \\
            &   &    *++[o][F-]{VA} \ar [lu]_{boire}^1  \\
            *++[o][F-]{VO}  \ar [ruu]^{boire}_1 \ar [r]^{empoisonner}_{0.1} \ar[rdd]^{vider}_{0.5} &*++[o][F-]{VE} \ar[ru]^{antidote}_1 \ar[rd]^{boire}_{0.1} \ar[dd]^{vider}_{0.5}&  \\
            &    & *++[o][F-]{Mort} \ar @(r,u) []_{\emptyset}^1 \\
             &     *++[o][F-]{VV} \ar [ru]^{\emptyset}_{0.5} &  \\
        }
    }
\end{align}

Il est à noter que cet exemple représente un exemple jouet, des choix subjectifs ont été effectués sur la représentation des états, des événements mais aussi concernant les valeurs attribuées par $expert\_judge_\tau$ qui correspondent au respect d'une norme sociale abstraite (et non statistique). On considère que les formules s'applique sur un exemple fourni sans le remettre en question.

\subsection{Définition des contextes causaux}
On réduit le champs de travail en ne travaillant que sur des contextes causaux partageant le même état de départ $S_0$ et une politique correcte relativement au LTS $<\mathds{S},\tau>$
\begin{itemize}
\item Contexte causal  $\chi_1 = (\pi_1,S_0,\tau)$ avec $\pi_1\in \Pi_\tau:$


    \begin{minipage}{0.45\textwidth}
        \begin{align*}
            &\pi_1(VO)= \{empoisonner\}, \quad \pi_1(VE)= \{boire\}, \quad \pi_1(mort)=\{\emptyset\} \quad \\ 
            &\pi_1(VV) = \{\emptyset\},\quad \pi_1(VA) = \{boire\},\quad \pi_1(vivant)=\{\emptyset\}
        \end{align*}
        On obtient:
    \begin{itemize}
        \item $\sigma^{\chi_1}= [VO,empoisonner,VE,boire, Mort]$
        \item $\sigma^{\chi_1}_{|state}=[VO,VE,Mort]$
        \item $\sigma^{\chi_1}_{|event}=[empoisonner,boire]$
    \end{itemize}
    \end{minipage}   
    \hfill
    \begin{minipage}{0.45\textwidth}
    \scalebox{0.8}{    \xymatrix@R=1em@C=5em{
            &    *++[o][F-]{Vivant} \ar @(r,u) []_{\emptyset}^1& \\
            &   &    *++[o][F-]{VA} \ar [lu]_{boire}^1  \\
            *++[o][F-]{VO}  \ar [ruu]^{boire}_1 \ar @[red][r]^{empoisonner}_{0.1} \ar[rdd]^{vider}_{0.5} &*++[o][F-]{VE} \ar[ru]^{antidote}_1 \ar@[red][rd]^{boire}_{0.1} \ar[dd]^{vider}_{0.5}&  \\
            &    & *++[o][F-]{Mort} \ar @(r,u) []_{\emptyset}^1 \\
             &     *++[o][F-]{VV} \ar [ru]^{\emptyset}_{0.5} &  \\
        }}

    \end{minipage}


    \item Contexte causal $\chi_2 = (\pi_2,S_0,\tau)$avec $\pi_2\in \Pi_\tau:$

    \begin{minipage}{0.45\textwidth}

    \begin{align*}
        &\pi_2(VO)= \{boire\}, \quad \pi_2(VE)= \{boire\}, \quad \pi_2(mort)=\{\emptyset\} \quad \\ 
        &\pi_2(VV) = \{\emptyset\},\quad \pi_2(VA) = \{boire\},\quad \pi_2(vivant)=\{\emptyset\}
    \end{align*}

    \begin{itemize}
        \item $\sigma^{\chi_2}= [VO,boire, Vivant]$
        \item $\sigma^{\chi_2}_{|state}=[VO,Vivant]$
        \item $\sigma^{\chi_2}_{|event}=[boire]$
    \end{itemize}
    \end{minipage}   
    \hfill
    \begin{minipage}{0.45\textwidth}
    \scalebox{0.8}{    \xymatrix@R=1em@C=5em{
            &    *++[o][F-]{Vivant} \ar @(r,u) []_{\emptyset}^1& \\
            &   &    *++[o][F-]{VA} \ar [lu]_{boire}^1  \\
            *++[o][F-]{VO}  \ar @[red] [ruu]^{boire}_1 \ar [r]^{empoisonner}_{0.1} \ar[rdd]^{vider}_{0.5} &*++[o][F-]{VE} \ar[ru]^{antidote}_1 \ar[rd]^{boire}_{0.1} \ar[dd]^{vider}_{0.5}&  \\
            &    & *++[o][F-]{Mort} \ar @(r,u) []_{\emptyset}^1 \\
             &     *++[o][F-]{VV} \ar [ru]^{\emptyset}_{0.5} &  \\
        }}

    \end{minipage}


    \item Contexte causal $\chi_3 = (\pi_3,S_0,\tau)$ avec $\pi_3\in \Pi_\tau:$

    \begin{minipage}{0.45\textwidth}
        \begin{align*}
            &\pi_3(VO)= \{vider\}, \quad \pi_3(VE)= \{boire\}, \quad \pi_3(mort)=\{\emptyset\} \quad \\ 
            &\pi_3(VV) = \{\emptyset\},\quad \pi_3(VA) = \{boire\},\quad \pi_3(vivant)=\{\emptyset\}
        \end{align*}
        \begin{itemize}
            \item $\sigma^{\chi_3}= [VO,vider, VV,\emptyset,Mort]$ 
            \item $\sigma^{\chi_3}_{|state}=[VO,VV,Mort]$
            \item $\sigma^{\chi_3}_{|event}=[vider,\emptyset]$
        \end{itemize}
        \end{minipage}   
        \hfill
        \begin{minipage}{0.45\textwidth}
        \scalebox{0.8}{    \xymatrix@R=1em@C=5em{
                &    *++[o][F-]{Vivant} \ar @(r,u) []_{\emptyset}^1& \\
                &   &    *++[o][F-]{VA} \ar [lu]_{boire}^1  \\
                *++[o][F-]{VO}  \ar  [ruu]^{boire}_1 \ar [r]^{empoisonner}_{0.1} \ar@[red][rdd]^{vider}_{0.5} &*++[o][F-]{VE} \ar[ru]^{antidote}_1 \ar[rd]^{boire}_{0.1} \ar[dd]^{vider}_{0.5}&  \\
                &    & *++[o][F-]{Mort} \ar @(r,u) []_{\emptyset}^1 \\
                 &     *++[o][F-]{VV} \ar @[red][ru]^{\emptyset}_{0.5} &  \\
            }}
    
        \end{minipage}

    \item Contexte causal 
    $\chi_4 = (\pi_4,S_0, \tau)$ avec $\pi_4\in \Pi_\tau:$

    \begin{minipage}{0.45\textwidth}
    \begin{align*}
        &\pi_4(VO)= \{empoisonner\}, \quad \pi_4(VE)= \{vider\}, \quad \pi_4(mort)=\{\emptyset\} \quad \\ 
        &\pi_4(VV) = \{\emptyset\},\quad \pi_4(VA) = \{boire\},\quad \pi_4(vivant)=\{\emptyset\}
    \end{align*}
    \begin{itemize}
        \item $\sigma^{\chi_4}= [VO,empoisonner,VE,vider,VV,\emptyset,Mort]$
        \item $\sigma^{\chi_4}_{|state}=[VO,VE,VV,Mort]$
        \item $\sigma^{\chi_4}_{|event}=[empoisonner,vider,\emptyset]$
    \end{itemize}        
    \end{minipage}
    \hfill
    \begin{minipage}{0.45\textwidth}
        \scalebox{0.8}{    \xymatrix@R=1em@C=5em{
                &    *++[o][F-]{Vivant} \ar @(r,u) []_{\emptyset}^1& \\
                &   &    *++[o][F-]{VA} \ar [lu]_{boire}^1  \\
                *++[o][F-]{VO}  \ar  [ruu]^{boire}_1 \ar@[red] [r]^{empoisonner}_{0.1} \ar[rdd]^{vider}_{0.5} &*++[o][F-]{VE} \ar[ru]^{antidote}_1 \ar[rd]^{boire}_{0.1} \ar@[red][dd]^{vider}_{0.5}&  \\
                &    & *++[o][F-]{Mort} \ar @(r,u) []_{\emptyset}^1 \\
                 &     *++[o][F-]{VV} \ar@[red][ru]^{\emptyset}_{0.5} &  \\
            }}
    \end{minipage}

    

    \item Contexte causal $\chi_5 = (\pi_5,S_0,\tau)$ avec $\pi_5\in \Pi_\tau:$


     \begin{minipage}{0.3\textwidth}
        \begin{align*}
            &\pi_5(VO)= \{empoisonner\}, \quad \pi_5(VE)= \{antidote\}, \quad \pi_5(mort)=\{\emptyset\} \quad \\ 
            &\pi_5(VV) = \{\emptyset\},\quad \pi_5(VA) = \{boire\},\quad \pi_5(vivant)=\{\emptyset\}
        \end{align*}
         \begin{itemize}
             \item $\sigma^{\chi_5}= [VO,empoisonner,VE,antidote,VA,boire,Vivant]$ 
             \item  $\sigma^{\chi_5}_{|state}=[VO,VE,VA,Vivant]$
             \item $\sigma^{\chi_5}_{|event}=[empoisonner,antidote,boire]$
         \end{itemize}         
     \end{minipage}
     \begin{minipage}{0.45\textwidth}
        \scalebox{0.8}{    \xymatrix@R=1em@C=5em{
                &    *++[o][F-]{Vivant} \ar @(r,u) []_{\emptyset}^1& \\
                &   &    *++[o][F-]{VA} \ar @[red][lu]_{boire}^1  \\
                *++[o][F-]{VO}  \ar  [ruu]^{boire}_1 \ar@[red] [r]^{empoisonner}_{0.1} \ar[rdd]^{vider}_{0.5} &*++[o][F-]{VE} \ar@[red][ru]^{antidote}_1 \ar[rd]^{boire}_{0.1} \ar[dd]^{vider}_{0.5}&  \\
                &    & *++[o][F-]{Mort} \ar @(r,u) []_{\emptyset}^1 \\
                 &     *++[o][F-]{VV} \ar[ru]^{\emptyset}_{0.5} &  \\
            }}
     \end{minipage}
\end{itemize}

\subsection{Définition des politiques optimales}

Politique optimale sans planification :
\begin{align*}
    & \pi^{opt1}_\tau(VO)=\{boire\}, \quad \pi^{opt1}_\tau(VE)=\{antidote\}, \quad \pi^{opt1}_\tau(VV)=\{\emptyset\}, \quad \pi^{opt1}_\tau(Mort)=\{\emptyset\} \\
    & \pi^{opt1}_\tau(VA)=\{boire\}, \quad \pi^{opt1}_\tau(Vivant)=\{\emptyset\} \quad
\end{align*}
Politique optimale avec planification:
\begin{align}\label{pol_opti}
    & \pi^{opt2}_\tau(VO)=\{boire\}, \quad \pi^{opt2}_\tau(VE)=\{antidote\}, \quad \pi^{opt2}_\tau(VV)=\{\emptyset\}, \quad \pi^{opt2}_\tau(Mort)=\{\emptyset\} \\
    & \pi^{opt2}_\tau(VA)=\{boire\}, \quad \pi^{opt2}_\tau(Vivant)=\{\emptyset\} \quad
\end{align}
On note $\forall S\in \mathds{S},\quad \pi^{opt1}_\tau(S) =\pi^{opt2}_\tau(S)$. On notera par la suite simplement $\pi^{opt}_\tau$.
\begin{figure}[h!]
    \centering
    \[
        \xymatrix@R=1em@C=5em{
            &    *++[o][F-]{Vivant} \ar@[red]  @(r,u) []_{\emptyset}^1& \\
            &   &    *++[o][F-]{VA} \ar@[red]  [lu]_{boire}^1  \\
            *++[o][F-]{VO}  \ar @[red] [ruu]^{boire}_1 \ar [r]^{empoisonner}_{0.1} \ar[rdd]^{vider}_{0.5} &*++[o][F-]{VE} \ar@[red] [ru]^{antidote}_1 \ar[rd]^{boire}_{0.1} \ar[dd]^{vider}_{0.5}&  \\
            &    & *++[o][F-]{Mort} \ar@[red]  @(r,u) []_{\emptyset}^1 \\
             &     *++[o][F-]{VV} \ar@[red] [ru]^{\emptyset}_{0.5} &  \\
        }    
    \]
    \captionsetup{width=0.8\textwidth}
    \caption{Politique optimale: les flèches rouges indiquent dans chaque état l'événement choisi par la politique optimale}
    \label{fig:placeholder}
\end{figure}

\subsection{Expérimentations}

En fonction de la configuration choisie (extraction d'une politique optimale, exploitation direct, évaluation de la normalité globale, sur les événements avec alternatives uniquement) et des hyper paramètres ($\lambda_1,\lambda_2, \gamma$...) on obtient différents résultats. Une connaissance à priori des contraintes que l'on souhaite appliquer lors de l'évaluation des elles permet de guider le classement. 

Dans les exemples suivant on considère, si ce n'est pas précisé, que le contexte causal observé est $\chi_1$, on fixe  $\lambda_1 =\lambda_2 =1$ et un coût maximal $\delta=2$ (avec cette valeur la validité ne dépends que de la différence entre les deux états finaux des traces de l'observation et du contrefactuel).

\paragraph{Priorité à des contrefactuels avec des scores d'événements consistants:}
On choisit pour la fonction de vraisemblance la formule \ref{exploit_direct} avec comme opérateur d'agrégation $agg_{geo}$ et on agrège les scores de tous les événements des scénarios. On souhaite dans cette configuration pénaliser des traces avec une variance de score importante.

On obtient les résultats suivant:

\begin{table}[h]
\centering
\begin{tabular}{|c|c|c|c|c|c|c|}
\hline
CF & classement & $d_{\text{obs}}$ & $v$ & $d_{\text{obs}} + v$ & trace & valide \\
\hline\hline
1 & $\chi_2$ & 0.5 & 0.0 & 0.5 &  $[VO,boire, Vivant]$ & \checkmark \\
\hline
3 & $\chi_3$ & 0.5 & 0.5 & 1 & $[VO,vider, VV,\emptyset,Mort]$ & \texttimes \\
4 & $\chi_4$ & 0.33 & 0.71 & 1.04 & $[VO,empoisonner,VE,vider,VV,\emptyset,Mort]$ &\texttimes \\
2 & $\chi_5$ & 0.33 & 0.54 & 0.87 & $[VO,empoisonner,VE,antidote,VA,boire,Vivant]$ & \checkmark \\
\hline
\end{tabular}
\caption{Résultats pour la priorité à la normativité}
\label{tab:prio_normativite}
\end{table}

On observe dans la table \ref{tab:prio_normativite} que le contexte contrefactuel avec le coût minimal est le contexte $\chi2$ sa distance à l'observation est de 0.5 et sa vraisemblance de 0 (l'unique événement $boire$ à un score de 1). Ce contrefactuel est suivi du contexte $\chi_5$, plus proche de l'observation (coût de 0.33) mais le scénario associé contient l'événement $empoisonner$ qui fait augmenter le coût. Les deux premiers contrefactuels $\chi_2$ et $\chi_5$ sont valides.



\paragraph{Interdiction de dépasser un seuil}

On utilise l'opérateur d'agrégation avec un seuil fixé à $\lambda = 0.2$ pour pénaliser les contrefactuels dont le scénario contient un empoisonnement. On fixe aussi un coût maximal de 0.5 pour le coût de vraisemblance. 

\begin{table}[h]
\centering
\begin{tabular}{|c|c|c|c|c|c|c|}
\hline
CF & classement & $d_{\text{obs}}$ & $v$ & $d_{\text{obs}} + v$ &  trace & valide \\
\hline\hline
1 & $\chi_2$ & 0.5 & 0.0 & 0.5 & $[VO,boire, Vivant]$ &  \checkmark \\
\hline
2 & $\chi_3$ & 0.5 & 0.0 & 0.5 & $[VO,vider, VV,\emptyset,Mort]$ & \texttimes \\
3 & $\chi_4$ & 0.33 & 1.0 & 1.33 & $[VO,empoisonner,VE,vider,VV,\emptyset,Mort]$ & \texttimes \\
4 & $\chi_5$ & 0.33 & 1.0 & 0.33 & $[VO,empoisonner,VE,antidote,VA,boire,Vivant]$ &  \texttimes \\
\hline
\end{tabular}
\caption{Résultats pour la priorité à la normativité}
\label{tab:res_seuil}
\end{table}

Dans cette configuration uniquement le contexte causal $\chi_2$ est valide. En rajoutant une contrainte sur le coût de vraisemblance on exclue automatiquement les contrefactuels qui possèdent un événement inférieur au seuil, c'est ce qu'on observe pour $\chi_5$. 

\paragraph{Normativité en début de scénario}
On utilise maintenant l'opérateur $agg_\gamma$, avec $\gamma = 0.25$, qui avantage les scénarios qui ont des scores élévées sur les événements en début de trace en considérant peu ceux de la fin de la trace.

\begin{table}[h]
\centering
\begin{tabular}{|c|c|c|c|c|c|c|}
\hline
CF& classement & $d_{\text{obs}}$ & $v$ & $d_{\text{obs}} + v$ & trace & valide \\
\hline\hline
1 & $\chi_2$ & 0.5 & 0.0 & 0.5 & $[VO,boire, Vivant]$ & \checkmark \\
\hline
2 & $\chi_3$ & 0.5 & 0.5 & 1.0 & $[VO,vider, VV,\emptyset,Mort]$ & \texttimes \\
4 & $\chi_4$ & 0.33 & 0.81 & 1.14 & $[VO,empoisonner,VE,vider,VV,\emptyset,Mort]$ & \texttimes \\
3 & $\chi_5$ & 0.33 & 0.69 & 1.02 & $[VO,empoisonner,VE,antidote,VA,boire,Vivant]$ & \checkmark \\
\hline
\end{tabular}
\caption{Priorité à la normativité en début de scénario}
\label{tab:res_gamma}
\end{table}

On observe avec la table \ref{tab:res_gamma} que le contrefactuel avec le plus faible coût corresponds à $\chi_2$ qui a la premier événement avec un score le plus élevé ($boire$ =1) suivit de $\chi_3$ qui a un score pour son premier événement de 0.5. Les deux derniers contextes causaux possède en premier événement $empoisonner$ ce qui vient fortement augmenter le coût. Le contexte $\chi_5$ a un coût plus faible que $\chi_4$ car, on baisse l'impact des événements en fin de scénario mais il reste présent ce qui permet une forme de compensation ($\chi_5$ utilise l'$antidote$).

\paragraph{Normativité en fin de scénario}

On utilise l'opérateur $gamma_{late}$ avec $\gamma =0.25$. On priorise la normativité en fin de trace, couplé au coût sur la distance à l'observation on cherche des contrefactuels qui ne diffèrent de l'observation qu'en fin de trace mais qui reste le plus proche possible au début.

\begin{table}[h]
\centering
\begin{tabular}{|c|c|c|c|c|c|c|}
\hline
CF & classement & $d_{\text{obs}}$ & $v$ & $d_{\text{obs}} + v$ & trace & valide \\
\hline\hline
2 & $\chi_2$ & 0.5 & 0.0 & 0.5 & $[VO,boire, Vivant]$ & \checkmark \\
4 &$\chi_3$ & 0.5 & 0.5 & 1.0 & $[VO,vider, VV,\emptyset,Mort]$ &  \texttimes \\
3 & $\chi_4$ & 0.33 & 0.81 & 0.85 & $[VO,empoisonner,VE,vider,VV,\emptyset,Mort]$ & \texttimes \\
\hline
1 & $\chi_5$ & 0.33 & 0.04 & 0.38 & $[VO,empoisonner,VE,antidote,VA,boire,Vivant]$ &  \checkmark \\
\hline\hline
\end{tabular}
\caption{Priorité à la normativité en fin de scénario}
\label{tab:res_gamma_late}
\end{table}

La table des résultats \ref{tab:res_gamma_late} montre que dans cette configuration c'est $\chi_5$ qui a le coût le plus faible, malgré l'empoisonnement on s'éloigne de l'observation le plus tard possible. On remarque que $\chi_3$ est le contexte causal le plus pénalisé avec cette configuration même si il n'y a pas d'empoisonnement dans ce dernier, cela est du au coût de $d\_obs$.

\paragraph{Normativité du contrefactuel avec distance à une politique optimale}
On choisit pour cette configuration de calculer la distance à la politique optimal définie par \ref{pol_opti}. Le choix est fait de ne considérer que les événements pour lesquelles une alternative existe.

\begin{table}[h]
\centering
\begin{tabular}{|c|c|c|c|c|c|c|}
\hline
CF& classement & $d_{\text{obs}}$ & $v$ & $d_{\text{obs}} + v$ & trace & valide \\
\hline
1 & $\chi_2$ & 0.5 & 0.0 & 0.5 & $[VO,boire, Vivant]$ & \checkmark \\
\hline
4 & $\chi_3$ & 0.5 & 1.0 & 1.5 & $[VO,vider, VV,\emptyset,Mort]$ & \texttimes \\
3 & $\chi_4$ & 0.33 & 1.0 & 1.33 & $[VO,empoisonner,VE,vider,VV,\emptyset,Mort]$ & \texttimes \\
2 & $\chi_5$ & 0.33 & 0.5 & 0.83 & $[VO,empoisonner,VE,antidote,VA,boire,Vivant]$ & \checkmark \\
\hline
\end{tabular}
\caption{Résultats distance à la politique optimale}
\label{tab:dist_pol}
\end{table}

Les résultats présents dans la table \ref{tab:dist_pol} indique qu'avec cette configuration, le contexte causal avec le plus faible coût est le contexte $\chi_2$ suivi de $\chi_5$ les deux contrefactuels valides. Concernant la fonction de vraisemblance: dans la trace associé à $\chi_2$ un seul événement est présent est un choix est possible, le choix fait est l'événement $boire$ qui corresponds à l'événement choisi par la politique optimale donc le coût associé est de 0. Concernant $\chi_5$, deux événements avec alternative sont présent dans la trace $empoisonner$ qui n'est pas l'événement indiqué par la politique optimal, et $antidote$ qui est similaire à la politique optimale. Sur les deux événements un seul est égal à la politique optimale donc un coût de 0.5 est appliqué. Pour $\chi_3$ et $\chi_4$ tous les choix possibles différent de la politique optimale un coût de 1 leur est attribué.
% \paragraph{Priorité à la proximité du contrefactuel à l'observation}








\newpage
\appendix

\section{Exemple PDDL}\label{ex_pddl}
Exemple d'instanciation de problème PDDL issus de \url{https://github.com/arii/planning_exercises/tree/master}. 

\subsection*{Domain Definition}
    \begin{lstlisting}[language=Lisp]
(define (domain air-cargo)
  (:requirements :strips)
  (:predicates 
    (In ?obj ?place)
    (At ?obj ?place)
    (Cargo ?obj)
    (Plane ?obj)
    (Airport ?obj))

  (:action LOAD
    :parameters (?c ?p ?a)
    :precondition (and (At ?c ?a) 
                       (At ?p ?a)
                       (Cargo ?c)
                       (Plane ?p)
                       (Airport ?a))
    :effect (and (not (At ?c ?a)) (In ?c ?p)))

  (:action UNLOAD
    :parameters (?c ?p ?a)
    :precondition (and (In ?c ?p)
                       (At ?p ?a)
                       (Cargo ?c)
                       (Plane ?p)
                       (Airport ?a))
    :effect (and (not (In ?c ?p)) (At ?c ?a)))

  (:action FLY
    :parameters (?p ?from ?to)
    :precondition (and (At ?p ?from)
                       (Plane ?p) 
                       (Airport ?from)
                       (Airport ?to))
    :effect (and (not (At ?p ?from)) (At ?p ?to))))
    \end{lstlisting}

\subsection*{Définition du problème}
    \begin{lstlisting}[language=Lisp]
; PDDL instance problem for the Air cargo transport
; Plane P1 is at SFO and plane P2 is at JFK.
; Move cargo C1 from SFO to JFK
; and cargo C2 from JFK to SFO.

(define (problem ac_p0)

  (:domain air-cargo)
  
  (:objects
    C1 C2   ; cargo objects
    P1 P2   ; plane objects
    SFO JFK ; airports
  )
	
  (:init
    ; types
    (Cargo C1) (Cargo C2)
    (Plane P1) (Plane P2)
    (Airport SFO) (Airport JFK)

    ; locations
    (At C1 SFO) (At C2 JFK)
    (At P1 SFO) (At P2 JFK))

  (:goal (and (At C1 JFK) (At C2 SFO))))
    \end{lstlisting}


\section{Exemple contexte causal non valide pour la fonction de coût}
Soit le LTS $\langle \mathds{S},\tau \rangle$ avec $$\mathds{S} = \{S_0,S_1,S_2,S_3\} \quad \tau = \{(S_0,a,S_1),(S_0,b,S_2),(S_2,c,S_3),(S_2,d,S_4),(S_1,\emptyset,S_1),(S_3,\emptyset,S_3),(S_4,\emptyset,S_4)\}$$:

\begin{figure}[h!]
    \centering
    \[
        \xymatrix @R=0.1em@C=5em {
        *++[o][F-]{S_0} \ar[r]^{a} \ar[rdd]^b & *++[o][F-]{S_1} \ar@(r,u)[]_{\emptyset}
        \\ & & *++[o][F-]{S_3} \ar@(r,u)[]_{\emptyset}
        \\ & *++[o][F-]{S_2} \ar[ru]^c \ar[rd]^d &
         \\ & & *++[o][F-]{S_4} \ar@(r,u)[]_{\emptyset}
        }    
    \]
    \caption{Exemple de LTS}
    \label{fig:placeholder}
\end{figure}

On définit:
$\chi_1 = (\pi_1,S_0,\tau)$, $(\pi_2,S_3,\tau)$ $(\pi_3,S_0,\tau)$et  avec: 
$$\pi_1(S_0) = a \quad \pi_1(S_2)=c \quad et \quad \forall i\in[1,3,4],\pi_1(S_i)= \emptyset$$
$$\pi_2(S_0) = a \quad \pi_2(S_2)=d \quad et \quad \forall i\in[1,3,4],\pi_2(S_i)= \emptyset$$
$$\pi_3(S_0) = e \quad \pi_3(S_2)=d \quad et \quad \forall i\in[1,3,4],\pi_3(S_i)= \emptyset$$



\subsection{Fonction $divergence$}\label{ex_div}

$divergence(\chi_1,\chi_2)= +\infty$ même si ces deux contextes causaux sont différents le coût est de $+\infty$, les politiques divergent uniquement en $S_2$ hors cet état est absent de $\sigma^{\chi_1}=[S_0,S_1]$. Les politiques sont égales sur la trace.

\subsection{Fonction $correct$}\label{ex_corr}
$correct(\chi_3) = +\infty$ la politique $\pi_3$ n'étant pas correcte relativement au LTS. $\nexists S' \textit{ tel que } (S_0,e,S')\in \tau$

\section{Construction d'une politique optimale avec planification et l'utilisation de l'algorithme A*}\label{astar}
\begin{figure}[h!]
    \centering
    \begin{subfigure}[b]{0.45\textwidth}
        \centering
        \[
            \xymatrix @R=0.1em@C=4em {
             & & *++[o][F-]{S_2}\ar@(r,u)[]_{\emptyset} 
             \\ *++[o][F-]{S_0} \ar[r]^{0.3}& *++[o][F-]{S_1} \ar[ru]^{0.7}  \ar[rd]^{0.4}   &
             \\ & &*++[o][F-]{S_3}\ar@(r,u)[]_{\emptyset} 
             }
        \]
        \caption{LTS original}
        \label{fig:diag-a}
    \end{subfigure}
    \hfill
    \begin{subfigure}[b]{0.45\textwidth}
        \centering
        \[
            \xymatrix @R=0.1em@C=4em {
             & & *++[o][F-]{S_2}\ar@(r,u)[]_{\emptyset} \ar@{-->} [rd]^{1} &
             \\ *++[o][F-]{S_0} \ar[r]^{0.3}& *++[o][F-]{S_1} \ar[ru]^{0.7}  \ar[rd]^{0.4}  &  & *++[o][F-]{temp}
             \\ & &*++[o][F-]{S_3}\ar@(r,u)[]_{\emptyset} \ar@{-->} [ru]^{1} &
             }
        \]
        \caption{LTS augmenté}
        \label{fig:diag-b}
    \end{subfigure}
    \caption{Construction du graphe pour l'application d'un algorithme de plus court chemin}
    \label{fig:lts astar}
\end{figure}
Les algorithmes de plus court chemin nécessitent l'indication d'un état de départ et d'arrivée. Dans notre cas on souhaite trouver le chemin qui minimise le coût sans contraindre à une destination fixe. Pour permettre cela on créer un LTS temporaire, à partir de notre LTS de travail, auquel on ajoute un état "poubelle" (voir Figure \ref{fig:lts astar}. On relie chaque état terminal du LTS original à ce nouvel état, et chaque transitions entre un état terminal et l'état possède un score égal (par exemple 1). L'ajout de cet état nous permet d'appliquer un algorithme de recherche du plus court chemin en fixant l'état d'arrivé à l'état ajouté.



\section{Vérification des propriétés pour les opérateurs}
Rappel:
Définition et propriétés d'un opérateur d'agrégation: $Agg: \underset{n\in \mathds{N}}{\cup}[0,1]^n\to [0,1]$
\begin{itemize}
    \item identité si $n=1 : Agg(u)=u$
    \item conditions aux limites: $Agg(0,\dots,0)=0$ et $Agg(1,\dots,1)=1$
    \item monotonie croissante: $Agg(u_1,\dots,u,\dots,u_n)\leq Agg(u_1,\dots,v,\dots,u_n)$ si $u \leq v$
\end{itemize}

\subsection{Opérateur géométrique}
Soit $\epsilon>0$
    \begin{equation}
            agg_{geo}(s_0,s_1,\dots,s_{n-1}) =\exp\left( \frac{1}{N} \sum_{i<n} \log(\max(s_i,\epsilon ))\right) 
    \end{equation}

\paragraph{Identité}
$$agg_{geo}(u) = \exp(\log(\max(u,\epsilon))) = \max(u,\epsilon)$$
Si $u<\epsilon$, $agg_{geo}(u)=\epsilon \neq u$. La propriété d'identité n'est pas respecté.

\paragraph{Conditions aux limites}
$$agg_{geo}(0,\dots,0) =\exp\left( \frac{1}{N} \sum_{i<n} \log(\epsilon)\right) =  \epsilon\neq 0$$

$$agg_{geo}(1,\dots,1) =\exp\left( \frac{1}{N} \sum_{i<n} \log(1)\right) =  1$$
Pour la limite en zéro la condition n'est pas respectée.

\paragraph{Monotonie croissante}
Soit $u \leq v$, $Agg(u_1,\dots,u,\dots,u_n)$ ne diffère de $ Agg(u_1,\dots,v,\dots,u_n)$ que par un seul terme de la somme.
On veut montrer que $\log(\max(u,\epsilon))\leq \log(\max(v,\epsilon))$ , hors $\max(u,\epsilon) \leq \max(v,\epsilon)$ et donc par croissance du log on obtient: $\log(\max(u,\epsilon))\leq \log(\max(v,\epsilon))$ . Donc la somme est plus petite et l'exponentielle aussi. La propriété de monotonie croissante est respectée.

\subsection{Opérateur gamma}
Soit $\gamma\in ]0,1]$, $\Gamma = \frac{1-\gamma^{N}}{1-\gamma}$ ($\Gamma =N $ si $\gamma=1$)
    \begin{equation}
        agg_{\gamma}(s_0,s_1,\dots,s_{n-1}) = \frac{ \sum_{i=0}^{N-1}\gamma^is_i}{\Gamma}
    \end{equation}
\paragraph{Identité}
Soit $u\in [0,1]$
$$agg_\gamma(u)= \frac{\gamma^0u}{\gamma^0}= u$$
La propriété d'identité est respectée.
\paragraph{Conditions aux limites}
$$agg_\gamma(0,\dots,0) = \frac{ \sum_{i=0}^{N-1}\gamma^i0}{\Gamma}=0$$

$$agg_\gamma(1,\dots,1) = \frac{ \sum_{i=0}^{N-1}\gamma^i1}{\Gamma}=1$$
Les conditions aux limites sont respectées.
\paragraph{Monotonie croissante}
Soit $u,v\in [0,1]$ tel que $u\leq v$ 
$$agg_{\gamma}(u_0,\dots,u,\dots,s_{n-1}) \leq agg_{\gamma}(u_0,\dots,v,\dots,s_{n-1})$$ car $\gamma^iu\leq \gamma^iv$
La propriété de monotonie croissante est respectée.

\subsection{Opérateur de seuil}

Soit $\lambda \in ]0,1]$

\begin{equation}
    agg_{seuil}(s_0,\dots,s_{n-1}) = \begin{cases}
        0 & si \quad \exists i<n \textit{ tel que } s_i < \lambda
        \\ 1 & sinon
    \end{cases}
\end{equation}

\paragraph{Identité}
Soit $u\in [0,1]$
$$agg_{seuil}=\begin{cases}
        0 & si \quad \exists i<n \textit{ tel que } u < \lambda
        \\ 1 & sinon
        \end{cases}$$
Si $u\neq 0$ ou $u\neq1$ la propriété d'identité n'est pas respectée.

\paragraph{Conditions aux limites}


$$agg_{seuil}(0,\dots,0) = 0 \quad agg_{seuil}(1,\dots,1) = 1 $$    
Les conditions aux limites sont respectées

\paragraph{Monotonie croissante}

Soit $u\leq v$, on veut montrer que: \begin{equation*}
    agg_{seuil}(u_0,\dots,u,\dots,u_{n-1}) \leq agg_{seuil}(u_0,\dots,v,\dots,u_{n-1})
\end{equation*}

\begin{itemize}
    \item Si $agg_{seuil}(u_0,\dots,v,\dots,u_{n-1})=1$ on a $agg_{seuil}(u_0,\dots,u,\dots,u_{n-1}) \leq agg_{seuil}(u_0,\dots,v,\dots,u_{n-1})$ 
    \item Si $agg_{seuil}(u_0,\dots,v,\dots,u_{n-1})=0$ alors il existe $i\in [0,n-1]$ tel que $u_i< \lambda$:
    \begin{itemize}
        \item Si $i$ est différents de l'indice de position de $u$ et $v$, alors ce même $u_i$ amène aussi à: \begin{equation*}
            agg_{seuil}(u_0,\dots,u,\dots,u_{n-1})=0
        \end{equation*} 
        \item Si $i$ corresponds à l'indice de position de $u$ et $v$ on a $u\leq v < \lambda$ et donc:\begin{equation*}
            agg_{seuil}(u_0,\dots,u,\dots,u_{n-1})=0
        \end{equation*} 
    \end{itemize}
    
\end{itemize}

La propriété de monotonie croissante est respectée.


\subsection{Synthèse résultats}
\begin{table}[h!]
    \centering
    \begin{tabular}{lccc}
        \hline
        Opérateur & Identité & Conditions aux limites & Monotonie croissante \\
        \hline
        $agg_{geo}$ & \texttimes & \texttimes\ (en 0) / \checkmark\ (en 1) & \checkmark \\
        $agg_{\gamma}$ & \checkmark & \checkmark & \checkmark \\
        $agg_{seuil}$ & \texttimes & \checkmark & \checkmark \\
        \hline
    \end{tabular}
    \caption{Synthèse des propriétés des opérateurs d'agrégation}
    \label{tab:synthese_agg}
\end{table}
\newpage
\printbibliography
\end{document}


